\documentclass{article}
\pagestyle{empty}

%% >>> preamble
\usepackage{xcolor}
\usepackage{listings}
\definecolor{codebg}{HTML}{F7F6F2}
\definecolor{codekw}{HTML}{0072B2}
\definecolor{codestr}{HTML}{009E73}
\definecolor{codecom}{HTML}{8A8577}
\lstdefinestyle{mystyle}{
  backgroundcolor=\color{codebg},
  basicstyle=\ttfamily\small,
  keywordstyle=\color{codekw}\bfseries,
  stringstyle=\color{codestr},
  commentstyle=\color{codecom}\itshape,
  numbers=left, numberstyle=\tiny\color{codecom}, numbersep=8pt,
  frame=single, rulecolor=\color{codebg!70!black},
  breaklines=true, showstringspaces=false, tabsize=4,
  xleftmargin=1.8em, framexleftmargin=1.5em, captionpos=b
}
\lstset{style=mystyle}
%% <<< preamble

\begin{document}
%% >>> snippet
\begin{lstlisting}[language=Python, caption={Fibonacci with memoisation.}]
from functools import lru_cache

@lru_cache(maxsize=None)
def fib(n: int) -> int:
    """Return the n-th Fibonacci number."""
    if n < 2:
        return n
    return fib(n - 1) + fib(n - 2)

print(fib(30))  # 832040
\end{lstlisting}
Inline code works too: \lstinline|fib(30)|.
%% <<< snippet
\end{document}
